% Unit 9 self-study -- "I do / You do" ladder for CED 9.1-9.11.
% Twelve ladders, forty-eight questions.
%
% Disjoint from u09-frq, which uses dH=+95.0/dS=+150, the Cu/Ag cell,
% diamond-to-graphite, NaCl dissolution and a +72/-105 coupling. This
% document uses the ammonia synthesis and CaCO3 decomposition for Gibbs,
% Fe/Cu and Ni/Ag cells, and different coupling data.
\documentclass{apchem-worksheet}

\apcunit{9}
\apcunittitle{Thermodynamics and Electrochemistry}
\apcdoctitle{Self-Study \textbullet\ CED 9.1--9.11, I~do / You~do}
\apcsubtitle{Twelve ladders \textbullet\ four YOUR TURN questions each \textbullet\ work alone, check after all four}

\begin{document}
\apctitleblock

\begin{apcnote}[How to use these notes --- read this first]
Each skill is a \textbf{ladder}: a worked example, then \textbf{four}
YOUR TURN questions --- same skill, new numbers, no help. Work all four
before checking anything.

\textbf{This unit is really one idea with two faces.} Thermodynamics asks
whether a reaction is favourable; electrochemistry measures that same
favourability as a voltage. The bridge is a single equation,
$\Delta G^\circ = -nFE^\circ$, and once you see that, the second half of
the unit is the first half with a voltmeter attached.

\textbf{The unit-conversion trap that costs the most marks.} $\Delta H$ is
tabulated in \textbf{kilojoules} and $\Delta S$ in \textbf{joules} per
kelvin. Mixing them in $\Delta G = \Delta H - T\Delta S$ produces an
answer wrong by a factor of a thousand. Convert first, every time.

$R = \SI{8.314}{\joule\per\mole\per\kelvin}$,
$F = \SI{96485}{\coulomb\per\mole}$.

\textbf{One scope note.} The CED excludes labelling an electrode as
positive or negative, since the convention differs between cell types.
Anode and cathode \emph{are} assessed, and the rule never changes:
oxidation at the anode, in both kinds of cell.
\end{apcnote}

% =====================================================================
\section*{\sffamily Ladder 1 \textbullet\ Predicting the sign of $\Delta S$}
\ced{9.1}

Entropy measures the number of ways energy and matter can be arranged.
More particles, more freedom, more entropy --- and \textbf{gases dominate}
every comparison they appear in.

\begin{workedexample}[I do: three processes]
Predict the sign of $\Delta S^\circ$ for:
(i) \ce{H2O(l) -> H2O(g)};
(ii) \ce{2NO2(g) -> N2O4(g)};
(iii) \ce{NaCl(s) -> Na+(aq) + Cl^-(aq)}.

\textbf{(i) Positive.} A liquid becomes a gas --- an enormous increase in
the freedom of the molecules. Phase changes to gas always raise entropy
sharply.

\textbf{(ii) Negative.} Two moles of gas become one. Fewer gas particles
means fewer arrangements.

\textbf{(iii) Positive.} Ions leave a highly ordered crystal lattice and
disperse through the solvent.

\textbf{The shortcut worth internalizing:} if the number of moles of
\emph{gas} changes, that decides the sign almost every time. Only when the
gas count is unchanged do you look at phases and particle counts.
\end{workedexample}

\begin{yourturn}[YOUR TURN 1 --- four questions]
\begin{enumerate}[leftmargin=*,itemsep=0.5em]
  \item Sign of $\Delta S$ for \ce{CaCO3(s) -> CaO(s) + CO2(g)}:
        \workspace[1.4cm]
  \item Sign for \ce{N2(g) + 3H2(g) -> 2NH3(g)}: \workspace[1.4cm]
  \item Sign for water freezing: \workspace[1.4cm]
  \item Sign for \ce{H2(g) + Cl2(g) -> 2HCl(g)}, and why it is small.
        \workspace[1.8cm]
\end{enumerate}
\selfcheck{(a) positive \quad (b) negative \quad (c) negative \quad
(d) near zero --- 2 moles of gas become 2}
\keyonly{\begin{apcanswer}(a) \textbf{Positive} --- a gas is produced where
there was none.
(b) \textbf{Negative} --- four moles of gas become two.
(c) \textbf{Negative} --- a disordered liquid becomes an ordered crystal.
(d) \textbf{Close to zero}, and slightly positive in practice. Two moles
of gas become two moles of gas, so the dominant factor is unchanged and
only smaller effects --- differences in molecular complexity and mass ---
remain.\end{apcanswer}}
\end{yourturn}

% =====================================================================
\section*{\sffamily Ladder 2 \textbullet\ Calculating $\Delta S^\circ$}
\ced{9.2}

\[ \Delta S^\circ = \sum S^\circ(\text{products}) -
   \sum S^\circ(\text{reactants}) \]
Unlike $\Delta H_f^\circ$, the standard entropy of an element is
\textbf{not} zero.

\begin{workedexample}[I do: products minus reactants]
Find $\Delta S^\circ$ for \ce{C(s) + O2(g) -> CO2(g)} from
$S^\circ$: \ce{C(s)} \num{5.7}, \ce{O2(g)} \num{205.0},
\ce{CO2(g)} \num{213.7} \si{\joule\per\mole\per\kelvin}.

\[ \Delta S^\circ = 213.7 - (5.7 + 205.0) = 213.7 - 210.7
   = \mathbf{+3.0~J/mol\,K} \]

\textbf{Why so small?} One mole of gas becomes one mole of gas. With the
gas count unchanged, the dominant term cancels and only a small residue
remains.

\textbf{Contrast with intuition:} you might expect burning to be strongly
entropy-increasing. It is not, in this case --- the disorder was already
there in the \ce{O2}.

\textbf{Do not set the element to zero.} \ce{C(s)} has
$S^\circ = 5.7$, not 0. Absolute entropies are zero only at
\SI{0}{\kelvin}, and that is a different statement from the
$\Delta H_f^\circ$ convention.
\end{workedexample}

\begin{yourturn}[YOUR TURN 2 --- four questions]
\begin{enumerate}[leftmargin=*,itemsep=0.5em]
  \item Products total \num{400}, reactants \num{560}
        \si{\joule\per\mole\per\kelvin}. Find $\Delta S^\circ$.
        \workspace[1.4cm]
  \item Is that consistent with a reaction that consumes gas?
        \workspace[1.6cm]
  \item Why is $S^\circ$ of an element not zero? \workspace[1.6cm]
  \item Which has the larger $S^\circ$, \ce{H2O(l)} or \ce{H2O(g)}?
        \workspace[1.4cm]
\end{enumerate}
\selfcheck{(a) $-160$ \si{\joule\per\mole\per\kelvin} \quad (b) yes \quad
(c) it has real disorder above \SI{0}{\kelvin} \quad (d) the gas}
\keyonly{\begin{apcanswer}(a) $400 - 560 = \mathbf{-160}$
\si{\joule\per\mole\per\kelvin}.
(b) \textbf{Yes.} A negative $\Delta S^\circ$ of that size is exactly what
consuming moles of gas produces.
(c) Because $S^\circ$ is an \textbf{absolute} entropy, not a change from a
reference. Any substance above \SI{0}{\kelvin} has real thermal motion and
therefore real entropy. This differs from $\Delta H_f^\circ$, which is
defined as zero for elements purely by convention.
(d) The \textbf{gas}, by a large margin --- its molecules are free to move
independently.\end{apcanswer}}
\end{yourturn}

% =====================================================================
\section*{\sffamily Ladder 3 \textbullet\ Gibbs free energy}
\ced{9.3}

\[ \Delta G = \Delta H - T\Delta S \]
Negative $\Delta G$ means thermodynamically favourable. Convert
$\Delta S$ to kilojoules first.

\begin{workedexample}[I do: the conversion that decides the answer]
For the ammonia synthesis,
$\Delta H^\circ = \SI{-92.2}{\kilo\joule\per\mole}$ and
$\Delta S^\circ = \SI{-198.7}{\joule\per\mole\per\kelvin}$. Find
$\Delta G^\circ$ at \SI{298}{\kelvin}.

\textbf{Convert:} $\Delta S^\circ =
\SI{-0.1987}{\kilo\joule\per\mole\per\kelvin}$.

\[ \Delta G^\circ = -92.2 - (298)(-0.1987) = -92.2 + 59.2
   = \mathbf{-33.0~kJ/mol} \]

Negative, so \textbf{favourable} at room temperature.

\textbf{Watch the double negative.} $\Delta S$ is negative, so
$-T\Delta S$ is \emph{positive} --- it works against the reaction. The
favourable enthalpy is what carries it.

\textbf{If you forget the conversion:} $-92.2 - (298)(-198.7)$ gives about
$+59\,100$, wrong in both magnitude and sign, and it would reverse your
conclusion entirely.
\end{workedexample}

\begin{yourturn}[YOUR TURN 3 --- four questions]
\begin{enumerate}[leftmargin=*,itemsep=0.5em]
  \item $\Delta H = \SI{-50.0}{\kilo\joule}$,
        $\Delta S = \SI{+100}{\joule\per\kelvin}$, $T = 298$~K. Find
        $\Delta G$. \workspace[1.8cm]
  \item Is that reaction favourable? \workspace[1.4cm]
  \item $\Delta H = \SI{+30.0}{\kilo\joule}$,
        $\Delta S = \SI{-50}{\joule\per\kelvin}$. Favourable at any
        temperature? \workspace[1.8cm]
  \item What does $\Delta G = 0$ mean? \workspace[1.6cm]
\end{enumerate}
\selfcheck{(a) $-79.8$~kJ \quad (b) yes \quad (c) never \quad
(d) at equilibrium}
\keyonly{\begin{apcanswer}(a) $\Delta G = -50.0 - (298)(0.100) =
-50.0 - 29.8 = \mathbf{-79.8}$~kJ.
(b) \textbf{Yes} --- negative $\Delta G$.
(c) \textbf{Never.} $\Delta H$ is positive (opposes) and $\Delta S$ is
negative, so $-T\Delta S$ is also positive (also opposes). Both terms work
against it at every temperature, so $\Delta G$ can never become negative.
(d) The system is at \textbf{equilibrium} --- neither direction is
favoured, and $K = 1$ under standard conditions.\end{apcanswer}}
\end{yourturn}

% =====================================================================
\section*{\sffamily Ladder 4 \textbullet\ Temperature dependence}
\ced{9.3}

The four sign combinations give four behaviours. Only two of them depend
on temperature.

\begin{workedexample}[I do: the four cases and a crossover]
Tabulate the cases, then find the crossover temperature for
\ce{CaCO3 -> CaO + CO2}, with
$\Delta H^\circ = \SI{+178.3}{\kilo\joule}$ and
$\Delta S^\circ = \SI{+160.6}{\joule\per\kelvin}$.

\textbf{The four cases.}
$\Delta H<0$, $\Delta S>0$: favourable at \emph{all} temperatures.
$\Delta H>0$, $\Delta S<0$: favourable at \emph{no} temperature.
$\Delta H<0$, $\Delta S<0$: favourable at \textbf{low} $T$.
$\Delta H>0$, $\Delta S>0$: favourable at \textbf{high} $T$ --- this case.

\textbf{Crossover, where $\Delta G^\circ = 0$:}
\[ T = \frac{\Delta H^\circ}{\Delta S^\circ} = \frac{178.3}{0.1606}
   = \mathbf{1110~K} \]

Above about \SI{1110}{\kelvin} the decomposition becomes favourable ---
which is why a lime kiln runs near \SI{1200}{\kelvin} and not at room
temperature.

\textbf{Check at \SI{298}{\kelvin}:}
$178.3 - (298)(0.1606) = +130$~kJ --- strongly unfavourable, as expected
for limestone sitting on a shelf.
\end{workedexample}

\begin{yourturn}[YOUR TURN 4 --- four questions]
\begin{enumerate}[leftmargin=*,itemsep=0.5em]
  \item $\Delta H<0$, $\Delta S>0$. At which temperatures is it
        favourable? \workspace[1.4cm]
  \item $\Delta H<0$, $\Delta S<0$. Which temperatures?
        \workspace[1.4cm]
  \item $\Delta H = \SI{+40.0}{\kilo\joule}$,
        $\Delta S = \SI{+80}{\joule\per\kelvin}$. Crossover
        temperature? \workspace[1.6cm]
  \item Why does raising $T$ help only when $\Delta S$ is positive?
        \workspace[1.8cm]
\end{enumerate}
\selfcheck{(a) all \quad (b) low only \quad (c) \SI{500}{\kelvin} \quad
(d) $-T\Delta S$ is negative only then}
\keyonly{\begin{apcanswer}(a) \textbf{All} temperatures --- both terms
favour it.
(b) \textbf{Low} temperatures only. The favourable enthalpy must outweigh
the unfavourable $-T\Delta S$, which grows with $T$.
(c) $T = 40.0/0.080 = \mathbf{500}$~K.
(d) Because the temperature appears only in the $-T\Delta S$ term. If
$\Delta S$ is positive, that term is negative and grows more negative as
$T$ rises, eventually overcoming a positive $\Delta H$. If $\Delta S$ is
negative, the term is positive and raising $T$ makes matters
\emph{worse}.\end{apcanswer}}
\end{yourturn}

% =====================================================================
\section*{\sffamily Ladder 5 \textbullet\ Thermodynamic versus kinetic control}
\ced{9.4}

$\Delta G$ says whether a reaction \emph{can} go and how far. It says
nothing about \emph{when}.

\begin{workedexample}[I do: favourable and yet stuck]
A mixture of petrol vapour and air has a strongly negative
$\Delta G$, yet sits indefinitely in a fuel tank. Explain, and say what a
spark does.

\textbf{Thermodynamically favourable:} combustion releases a great deal of
energy and increases entropy. $\Delta G$ is large and negative.

\textbf{Kinetically blocked:} the reaction must begin by breaking strong
\ce{C-H}, \ce{C-C} and \ce{O=O} bonds --- a large activation energy. At
room temperature essentially no molecules have that much energy, so the
rate is effectively zero.

\textbf{What a spark supplies:} enough energy for a small number of
molecules to cross the barrier. Once started, the reaction is exothermic
and supplies the energy to keep itself going.

\textbf{The general principle:} $\Delta G$ describes the destination,
$E_a$ the journey. A large negative $\Delta G$ with a large $E_a$ is a
reaction that is strongly favourable and completely stuck --- which is
fortunate, since it is why fuels can be stored at all.
\end{workedexample}

\begin{yourturn}[YOUR TURN 5 --- four questions]
\begin{enumerate}[leftmargin=*,itemsep=0.5em]
  \item Does a negative $\Delta G$ guarantee a fast reaction?
        \workspace[1.4cm]
  \item What quantity controls the rate? \workspace[1.4cm]
  \item Diamond to graphite has $\Delta G^\circ$ negative. Why do diamonds
        persist? \workspace[1.8cm]
  \item Name two ways to speed a favourable but slow reaction.
        \workspace[1.6cm]
\end{enumerate}
\selfcheck{(a) no \quad (b) activation energy \quad (c) enormous $E_a$
\quad (d) raise $T$; add a catalyst}
\keyonly{\begin{apcanswer}(a) \textbf{No.} It guarantees only that the
reaction is thermodynamically favourable.
(b) The \textbf{activation energy} (together with temperature), not
$\Delta G$.
(c) Rearranging diamond's covalent network requires breaking very strong
\ce{C-C} bonds --- an enormous activation energy. At room temperature the
rate is indistinguishable from zero, so diamonds are said to be
\textbf{kinetically stable} even though graphite is the thermodynamically
favoured form.
(d) \textbf{Raise the temperature} (more molecules exceed $E_a$) and
\textbf{add a catalyst} (lowers $E_a$ by offering a new pathway). Neither
changes $\Delta G$.\end{apcanswer}}
\end{yourturn}

% =====================================================================
\section*{\sffamily Ladder 6 \textbullet\ Free energy and the equilibrium constant}
\ced{9.5}

\[ \Delta G^\circ = -RT\ln K \]
$\Delta G^\circ$ must be in \textbf{joules} to match $R$.

\begin{workedexample}[I do: both directions]
Find $K$ at \SI{298}{\kelvin} when
$\Delta G^\circ = \SI{-20.0}{\kilo\joule\per\mole}$. Then state $K$ when
$\Delta G^\circ = 0$.

\[ \ln K = -\frac{\Delta G^\circ}{RT}
   = \frac{20\,000}{(8.314)(298)} = \num{8.073} \]
\[ K = e^{8.073} = \mathbf{3.2\times10^{3}} \]

Large $K$, as a negative $\Delta G^\circ$ requires.

\textbf{When $\Delta G^\circ = 0$:} $\ln K = 0$, so $\mathbf{K = 1}$ ---
neither side favoured under standard conditions.

\textbf{The three-way correspondence to memorize:}
$\Delta G^\circ < 0 \Leftrightarrow K > 1 \Leftrightarrow$ products
favoured;
$\Delta G^\circ = 0 \Leftrightarrow K = 1$;
$\Delta G^\circ > 0 \Leftrightarrow K < 1 \Leftrightarrow$ reactants
favoured.

\textbf{The unit trap again:} leaving $\Delta G^\circ$ in kilojoules gives
$\ln K = 0.0081$ and $K \approx 1.008$ --- a wildly wrong answer that
still looks like a number.
\end{workedexample}

\begin{yourturn}[YOUR TURN 6 --- four questions]
\begin{enumerate}[leftmargin=*,itemsep=0.5em]
  \item $\Delta G^\circ = \SI{+15.0}{\kilo\joule\per\mole}$ at
        \SI{298}{\kelvin}. Find $K$. \workspace[1.8cm]
  \item Is $K$ greater or less than 1 when $\Delta G^\circ$ is negative?
        \workspace[1.4cm]
  \item What is $K$ when $\Delta G^\circ = 0$? \workspace[1.4cm]
  \item Why must $\Delta G^\circ$ be converted to joules?
        \workspace[1.6cm]
\end{enumerate}
\selfcheck{(a) \num{2.4e-3} \quad (b) greater than 1 \quad (c) 1 \quad
(d) to match $R$ in \si{\joule\per\mole\per\kelvin}}
\keyonly{\begin{apcanswer}(a) $\ln K = -15\,000/[(8.314)(298)] =
-6.054$; $K = e^{-6.054} = \mathbf{2.4\times10^{-3}}$.
(b) \textbf{Greater than 1} --- products are favoured.
(c) $K = \mathbf{1}$.
(d) Because $R = \SI{8.314}{\joule\per\mole\per\kelvin}$ carries joules.
Mixing kilojoules with it leaves the exponent a thousand times too small,
and $K$ comes out close to 1 no matter what the real answer
was.\end{apcanswer}}
\end{yourturn}

% =====================================================================
\section*{\sffamily Ladder 7 \textbullet\ Free energy of dissolution}
\ced{9.6}

Dissolving is a competition: breaking the lattice costs energy, hydration
releases it, and entropy almost always favours the process.

\begin{workedexample}[I do: endothermic yet spontaneous]
Ammonium nitrate dissolves readily although the solution gets markedly
colder. Explain.

\textbf{$\Delta H$ is positive.} The energy needed to pull the lattice
apart exceeds what hydration releases, so the process absorbs energy from
the water --- which is why it cools.

\textbf{$\Delta S$ is strongly positive.} Ions locked in an ordered
crystal become dispersed and mobile throughout the solvent --- a large
increase in the number of available arrangements.

\textbf{$\Delta G = \Delta H - T\Delta S$ is negative} because the
$T\Delta S$ term outweighs the positive $\Delta H$ at room temperature.
The process is entropy-driven.

\textbf{The general lesson:} spontaneity is not the same as
exothermicity. An endothermic process proceeds whenever the entropy gain
is large enough --- and for dissolving an ionic solid, it usually is.
\end{workedexample}

\begin{yourturn}[YOUR TURN 7 --- four questions]
\begin{enumerate}[leftmargin=*,itemsep=0.5em]
  \item Sign of $\Delta S$ when an ionic solid dissolves?
        \workspace[1.4cm]
  \item What drives an endothermic dissolution? \workspace[1.4cm]
  \item If a salt is insoluble despite favourable entropy, what must be
        true of $\Delta H$? \workspace[1.8cm]
  \item Would cooling the water make \ce{NH4NO3} more or less soluble?
        \workspace[1.8cm]
\end{enumerate}
\selfcheck{(a) positive \quad (b) the entropy increase \quad
(c) strongly positive \quad (d) less soluble}
\keyonly{\begin{apcanswer}(a) \textbf{Positive} --- an ordered lattice
becomes dispersed ions.
(b) The \textbf{entropy increase}: $T\Delta S$ outweighs the positive
$\Delta H$.
(c) It must be \textbf{strongly positive} --- large enough that the
endothermic cost of separating the lattice exceeds $T\Delta S$, leaving
$\Delta G$ positive.
(d) \textbf{Less soluble.} Lowering $T$ shrinks the favourable
$T\Delta S$ term while the unfavourable positive $\Delta H$ is unchanged,
so $\Delta G$ becomes less negative --- and eventually
positive.\end{apcanswer}}
\end{yourturn}

% =====================================================================
\section*{\sffamily Ladder 8 \textbullet\ Coupled reactions}
\ced{9.7}

Free energies add, exactly as enthalpies do in Hess's law. An unfavourable
reaction can be driven by pairing it with a strongly favourable one that
shares a species.

\begin{workedexample}[I do: driving an unfavourable step]
Reaction A has $\Delta G^\circ = \SI{+60}{\kilo\joule\per\mole}$;
reaction B, sharing a common species, has
$\Delta G^\circ = \SI{-95}{\kilo\joule\per\mole}$. Find the coupled value.

\[ \Delta G^\circ_{\text{total}} = (+60) + (-95) = \mathbf{-35~kJ/mol} \]

Negative, so the coupled process \textbf{is} favourable even though A
alone is not.

\textbf{The essential condition: a shared species.} The product of one
reaction must be consumed by the other. Without that link they are two
independent reactions that happen to be in the same beaker, and the
unfavourable one still will not proceed.

\textbf{Where this matters:} metal extraction couples an unfavourable
oxide decomposition to the favourable oxidation of carbon. In biology,
ATP hydrolysis is coupled to thousands of otherwise unfavourable
reactions --- it is how cells do chemistry that thermodynamics would
otherwise forbid.
\end{workedexample}

\begin{yourturn}[YOUR TURN 8 --- four questions]
\begin{enumerate}[leftmargin=*,itemsep=0.5em]
  \item $+45$ and $-70$ \si{\kilo\joule\per\mole}. Coupled value, and is
        it favourable? \workspace[1.6cm]
  \item $+80$ and $-55$. Coupled value, and is it favourable?
        \workspace[1.6cm]
  \item What condition must the two reactions satisfy?
        \workspace[1.6cm]
  \item Why is coupling analogous to Hess's law? \workspace[1.6cm]
\end{enumerate}
\selfcheck{(a) $-25$, favourable \quad (b) $+25$, not favourable \quad
(c) a shared species \quad (d) both are state functions and add}
\keyonly{\begin{apcanswer}(a) $\mathbf{-25}$~kJ/mol --- \textbf{favourable}.
(b) $\mathbf{+25}$~kJ/mol --- \textbf{not favourable}. The favourable
reaction is not strong enough to carry the other.
(c) They must share a \textbf{common species}, so that the product of one
is consumed by the other and they proceed as a single combined process.
(d) Because $G$, like $H$, is a \textbf{state function}: its change
depends only on the initial and final states. Values for steps that add to
an overall process therefore sum to the overall value, whichever route is
taken.\end{apcanswer}}
\end{yourturn}

% =====================================================================
\section*{\sffamily Ladder 9 \textbullet\ Galvanic cells}
\ced{9.8}

Oxidation at the anode, reduction at the cathode --- in both cell types.
Electrons flow through the wire from anode to cathode.

\begin{workedexample}[I do: assembling a cell]
Build a cell from iron ($E^\circ = -0.44$~V) and copper
($E^\circ = +0.34$~V). Give the half-reactions, the anode, $E^\circ$, and
the directions of electron and anion flow.

\textbf{The more positive potential is reduced}, so copper is the
cathode and iron is oxidized at the anode.

\ce{Fe(s) -> Fe^2+(aq) + 2e^-} (anode) \quad
\ce{Cu^2+(aq) + 2e^- -> Cu(s)} (cathode)

\[ E^\circ = E^\circ_{\text{cathode}} - E^\circ_{\text{anode}}
   = 0.34 - (-0.44) = \mathbf{+0.78~V} \]

\textbf{Electrons} flow through the external wire from the \textbf{iron}
to the \textbf{copper}.

\textbf{Anions} in the salt bridge migrate toward the \textbf{anode},
balancing the positive charge building up as \ce{Fe^2+} enters that
solution.

\textbf{Positive $E^\circ$ confirms it works} --- a galvanic cell always
has a positive cell potential, which is the same statement as
$\Delta G^\circ$ being negative.
\end{workedexample}

\begin{yourturn}[YOUR TURN 9 --- four questions]
\begin{enumerate}[leftmargin=*,itemsep=0.5em]
  \item Nickel ($-0.25$~V) and silver ($+0.80$~V): which is the anode, and
        what is $E^\circ$? \workspace[1.8cm]
  \item Which way do electrons flow in the external circuit?
        \workspace[1.4cm]
  \item Which way do anions move in the salt bridge?
        \workspace[1.4cm]
  \item What happens at the anode in an \emph{electrolytic} cell?
        \workspace[1.6cm]
\end{enumerate}
\selfcheck{(a) nickel; $+1.05$~V \quad (b) anode to cathode \quad
(c) toward the anode \quad (d) oxidation --- same as galvanic}
\keyonly{\begin{apcanswer}(a) \textbf{Nickel} is the anode (less positive
potential, so it is oxidized);
$E^\circ = 0.80 - (-0.25) = \mathbf{+1.05}$~V.
(b) From the \textbf{anode to the cathode} through the wire.
(c) Toward the \textbf{anode}, to balance the positive charge accumulating
as metal ions enter that solution.
(d) \textbf{Oxidation} --- exactly as in a galvanic cell. That never
changes. (What does differ between the cell types is the $+$/$-$
labelling, which is why the CED excludes
it.)\end{apcanswer}}
\end{yourturn}

% =====================================================================
\section*{\sffamily Ladder 10 \textbullet\ Cell potential and free energy}
\ced{9.9}

\[ \Delta G^\circ = -nFE^\circ \]
$n$ is the number of electrons transferred in the balanced reaction.
$E^\circ$ is \emph{intensive} --- never multiply it.

\begin{workedexample}[I do: the bridge equation]
For the iron--copper cell above ($E^\circ = +0.78$~V, $n = 2$), find
$\Delta G^\circ$.

\[ \Delta G^\circ = -nFE^\circ = -(2)(96485)(0.78)
   = \num{-1.5e5}~\text{J} = \mathbf{-151~kJ} \]

Negative, as it must be for a cell that delivers current on its own.

\textbf{The intensive property that trips everyone.} If you double the
whole reaction, $n$ doubles and $\Delta G^\circ$ doubles --- but
$E^\circ$ does \textbf{not}. Potential is energy per unit charge, so
scaling the reaction scales both energy and charge together and leaves the
ratio alone. A cell's voltage does not depend on its size.

\textbf{The correspondence:} $E^\circ > 0 \Leftrightarrow
\Delta G^\circ < 0 \Leftrightarrow K > 1$. Three ways of saying the same
thing.
\end{workedexample}

\begin{yourturn}[YOUR TURN 10 --- four questions]
\begin{enumerate}[leftmargin=*,itemsep=0.5em]
  \item $E^\circ = +1.05$~V with $n = 2$. Find $\Delta G^\circ$.
        \workspace[1.8cm]
  \item If a reaction is doubled, what happens to $E^\circ$ and to
        $\Delta G^\circ$? \workspace[1.8cm]
  \item What is the sign of $\Delta G^\circ$ for an electrolytic cell?
        \workspace[1.4cm]
  \item A cell has $E^\circ = -0.30$~V. What does that tell you?
        \workspace[1.8cm]
\end{enumerate}
\selfcheck{(a) $-203$~kJ \quad (b) $E^\circ$ unchanged, $\Delta G^\circ$
doubles \quad (c) positive \quad (d) the reverse reaction is the
favourable one}
\keyonly{\begin{apcanswer}(a) $-(2)(96485)(1.05) = \num{-2.03e5}$~J
$= \mathbf{-203}$~kJ.
(b) $E^\circ$ is \textbf{unchanged} --- it is intensive.
$\Delta G^\circ$ \textbf{doubles}, because $n$ doubles.
(c) \textbf{Positive.} An electrolytic cell drives a reaction that is not
thermodynamically favourable, which is why it needs an external power
supply.
(d) That the reaction as written is \textbf{not} favourable --- its
$\Delta G^\circ$ is positive. The \emph{reverse} reaction is the
spontaneous one, with $E^\circ = +0.30$~V.\end{apcanswer}}
\end{yourturn}

% =====================================================================
\section*{\sffamily Ladder 11 \textbullet\ The Nernst equation}
\ced{9.10}

\[ E = E^\circ - \frac{0.0592}{n}\log Q \qquad
   (\SI{25}{\celsius}) \]
Non-standard concentrations shift the potential. As the cell runs, $Q$
rises and $E$ falls.

\begin{workedexample}[I do: which way does it move?]
A Zn/Cu cell has $E^\circ = 1.10$~V and $n = 2$. Find $E$ when
$Q = 0.010$, and when $Q = 100$.

\textbf{$Q = 0.010$:}
\[ E = 1.10 - \frac{0.0592}{2}\log(0.010) = 1.10 - (0.0296)(-2)
   = \mathbf{1.16~V} \]
\emph{Above} standard --- $Q < 1$ means reactants are plentiful, so the
driving force is greater.

\textbf{$Q = 100$:}
\[ E = 1.10 - (0.0296)(2) = \mathbf{1.04~V} \]
\emph{Below} standard --- products have accumulated.

\textbf{What happens as the cell discharges:} products build up, $Q$ rises,
and $E$ falls steadily. A ``dead'' battery has reached equilibrium:
$Q = K$ and $E = \mathbf{0}$.

\textbf{Note carefully: $E = 0$, not $E^\circ = 0$.} The standard
potential is a constant of the chemistry and never changes; it is the
actual potential that decays to zero.
\end{workedexample}

\begin{yourturn}[YOUR TURN 11 --- four questions]
\begin{enumerate}[leftmargin=*,itemsep=0.5em]
  \item Does $E$ rise or fall as a cell discharges?
        \workspace[1.4cm]
  \item What is $E$ for a dead battery? \workspace[1.4cm]
  \item $E^\circ = 0.80$~V, $n = 1$, $Q = 10$. Find $E$.
        \workspace[1.8cm]
  \item Why does $E^\circ$ stay constant while $E$ changes?
        \workspace[1.8cm]
\end{enumerate}
\selfcheck{(a) falls \quad (b) zero \quad (c) \num{0.74}~V \quad
(d) $E^\circ$ is defined at standard conditions}
\keyonly{\begin{apcanswer}(a) It \textbf{falls}, as products accumulate
and $Q$ rises.
(b) $E = \mathbf{0}$ --- the cell has reached equilibrium, $Q = K$, and no
further net reaction occurs.
(c) $E = 0.80 - (0.0592/1)\log(10) = 0.80 - 0.0592 =
\mathbf{0.74}$~V.
(d) Because $E^\circ$ is \emph{defined} at standard conditions --- all
solutes at \SI{1}{\Molar}, gases at \SI{1}{\atm}. It is a fixed property
of the cell chemistry. $E$ is the potential under whatever conditions
actually exist, so it changes as those conditions
change.\end{apcanswer}}
\end{yourturn}

% =====================================================================
\section*{\sffamily Ladder 12 \textbullet\ Electrolysis and Faraday's law}
\ced{9.11}

Charge is current times time; moles of electrons is charge over $F$; then
the half-reaction's electron count converts to moles of substance.

\begin{workedexample}[I do: the four-step chain]
Copper is plated from \ce{Cu^2+} at \SI{2.50}{\ampere} for
\SI{60.0}{\minute}. Find the mass deposited.

\textbf{Time to seconds:} $60.0 \times 60 = \SI{3600}{\second}$.

\textbf{Charge:} $Q = It = (2.50)(3600) = \SI{9000}{\coulomb}$.

\textbf{Moles of electrons:}
$9000/96485 = \num{0.0933}$~mol.

\textbf{Through the half-reaction.} \ce{Cu^2+ + 2e^- -> Cu} needs
\textbf{two} electrons per atom:
\[ n(\ce{Cu}) = \frac{0.09328}{2} = \num{0.04664}~\text{mol} \]
\[ m = 0.04664 \times 63.55 = \mathbf{2.96~g} \]

\textbf{The step that is missed most often} is dividing by 2. Silver
(\ce{Ag+ + e^- -> Ag}) needs only one electron, so the same charge
deposits twice as many moles of silver as of copper.
\end{workedexample}

\begin{yourturn}[YOUR TURN 12 --- four questions]
\begin{enumerate}[leftmargin=*,itemsep=0.5em]
  \item Charge delivered by \SI{1.00}{\ampere} for \SI{30.0}{\minute}:
        \workspace[1.4cm]
  \item Moles of electrons in that charge: \workspace[1.6cm]
  \item Mass of silver deposited by it
        ($M = \SI{107.87}{\gram\per\mole}$): \workspace[1.8cm]
  \item Same charge through \ce{Al^3+}: more or fewer moles of metal than
        silver, and by what factor? \workspace[1.8cm]
\end{enumerate}
\selfcheck{(a) \SI{1800}{\coulomb} \quad (b) \num{0.0187}~mol \quad
(c) \SI{2.01}{\gram} \quad (d) one third as many}
\keyonly{\begin{apcanswer}(a) $Q = (1.00)(1800) =
\mathbf{1800}$~C.
(b) $1800/96485 = \mathbf{0.0187}$~mol.
(c) \ce{Ag+ + e^- -> Ag} is a one-electron reduction, so
$n(\ce{Ag}) = \num{0.01866}$~mol and
$m = 0.01866 \times 107.87 = \mathbf{2.01}$~g.
(d) \textbf{One third as many.} \ce{Al^3+ + 3e^- -> Al} requires three
electrons per atom, so the same supply of electrons produces a third of
the moles of metal.\end{apcanswer}}
\end{yourturn}

% =====================================================================
\section*{\sffamily Where you stand}

Tick a ladder only if all four YOUR TURN questions were right first time.

\renewcommand{\arraystretch}{1.30}
\noindent\begin{tabular}{@{}c p{6.0cm} c >{\raggedright\arraybackslash}p{4.9cm}@{}}
\toprule
\textbf{First try?} & \textbf{Skill} & \textbf{Ladder} &
  \textbf{If not, re-read\ldots}\\
\midrule
$\square$ & sign of $\Delta S$ & 1 & count moles of gas \\
$\square$ & calculating $\Delta S^\circ$ & 2 & elements are not zero \\
$\square$ & Gibbs free energy & 3 & convert $\Delta S$ to kJ \\
$\square$ & temperature dependence & 4 & the four sign cases \\
$\square$ & kinetic vs thermodynamic & 5 & how far vs how fast \\
$\square$ & $\Delta G^\circ$ and $K$ & 6 & joules, to match $R$ \\
$\square$ & dissolution & 7 & entropy can drive it \\
$\square$ & coupled reactions & 8 & needs a shared species \\
$\square$ & galvanic cells & 9 & oxidation at the anode \\
$\square$ & $\Delta G^\circ = -nFE^\circ$ & 10 & never scale $E^\circ$ \\
$\square$ & the Nernst equation & 11 & $E \to 0$, not $E^\circ$ \\
$\square$ & Faraday's law & 12 & divide by the electron count \\
\bottomrule
\end{tabular}

\begin{apcnote}[Scoring yourself honestly]
12/12: move on to the unit worksheets and the free-response set.
9--11: solid --- redo the missed ladders tomorrow from a blank page.
8 or fewer: the recurring root causes here are exactly three ---
(1) mixing kilojoules and joules in $\Delta G = \Delta H - T\Delta S$ or
in $\Delta G^\circ = -RT\ln K$, (2) multiplying $E^\circ$ when the
reaction is scaled, and (3) forgetting to divide by the number of
electrons in a Faraday calculation. Fix those three and most of these
ladders fall together.
\end{apcnote}

\end{document}
